\documentclass{article}
\usepackage{graphicx}
\usepackage{amsmath}
\usepackage{amsfonts}
\usepackage{bookmark}
\usepackage{booktabs}
\usepackage{exercises}
\graphicspath{ {./images/9.6} }
\usepackage{indentfirst}
\usepackage{amsthm}
\usepackage{yhmath}
\usepackage{geometry}
\usepackage{placeins}
\usepackage{csquotes}
\usepackage{pgfplots}
\newcommand*{\oldneg}{\mathord{\sim}}
\newcommand\showdiv[1]{\overline{\smash{)}#1}}
\newcommand\Mydiv[1]{%
\smash{\raise.3ex\hbox{$\big)$}}\mkern-8mu
        \overline{\enspace\strut#1}}
\setlength{\tabcolsep}{5pt}
\theoremstyle{definition}
\newtheorem{definition}{Definition}

\setlength{\parindent}{.5in}
\setlength{\parskip}{12pt}
\renewcommand{\arraystretch}{2}
\geometry{margin=1in}
\newcommand{\?}{\stackrel{?}{=}}
\pgfplotsset{width=7cm,compat=1.18}\usepgfplotslibrary{statistics}

\title{9.6 Answers}
\begin{document}
\maketitle

\section*{Derivation of Heron's Formula}
First, let us note that $\sin^2\theta+\cos^2\theta=1$. For consider an angle $\theta$ intersecting the unit circle at $(x,y)$. Then we know that $\cos\theta=x$ and 
$\sin\theta=y$. But by the Pythagorean theorem, $x^2+y^2=1$; hence $\sin^2\theta+\cos^2\theta=1$. This also means that
\[ \textrm{(1) }\sin\theta=\sqrt{1-\cos^2\theta} \] Now consider a triangle with side lengths $a,b,c$. By the law of cosines, we know that
\[ c^2=a^2+b^2-2ab\cos C \]
\[ 2ab\cos C=a^2+b^2-c^2 \]
\[ \textrm{(2) }\cos C=\dfrac{a^2+b^2-c^2}{2ab} \]
Also, we know that the area of the triangle is $A=\frac{1}{2}ab\sin C$. Thus:
\[ A=\frac{1}{2}ab\sin C \]
\[ =\frac{1}{2}ab\sqrt{1-\cos^2 C}\textrm{ (By 1)} \]
\[ =\frac{1}{2}ab\sqrt{1-\left(\dfrac{a^2+b^2-c^2}{2ab}\right)^2}\textrm{ (By 2)} \]
\[ =\frac{1}{2}ab\sqrt{1-\dfrac{(a^2+b^2-c^2)^2}{4a^2b^2}} \]
\[ =\frac{1}{2}ab\sqrt{1-\dfrac{(a^2+b^2-c^2)a^2+(a^2+b^2-c^2)b^2-(a^2+b^2-c^2)c^2}{4a^2b^2}} \]
\[ =\frac{1}{2}ab\sqrt{1-\dfrac{a^4+a^2b^2-a^2c^2+a^2b^2+b^4-b^2c^2-a^2c^2-b^2c^2+c^4}{4a^2b^2}} \]
\[ =\frac{1}{2}ab\sqrt{1-\dfrac{a^4+b^4+c^4+2a^2b^2-2a^2c^2-2b^2c^2}{4a^2b^2}} \]
\[ =\frac{1}{2}ab\sqrt{\dfrac{4a^2b^2-(a^4+b^4+c^4+2a^2b^2-2a^2c^2-2b^2c^2)}{4a^2b^2}} \]
\[ =\frac{1}{2}ab\sqrt{\dfrac{2a^2b^2-a^4-b^4-c^4+2a^2c^2+2b^2c^2}{4a^2b^2}} \]
\[ =\frac{1}{2}ab\dfrac{\sqrt{2a^2b^2-a^4-b^4-c^4+2a^2c^2+2b^2c^2}}{2ab} \]
\[ \dfrac{\sqrt{2a^2b^2-a^4-b^4-c^4+2a^2c^2+2b^2c^2}}{4} \]
\[ \sqrt{\dfrac{2a^2b^2-a^4-b^4-c^4+2a^2c^2+2b^2c^2}{16}} \]
Now, one can see that
\[ \left(a+b+c\right)\left(-a+b+c\right)\left(a-b+c\right)\left(a+b-c\right)=2a^2b^2-a^4-b^4-c^4+2a^2c^2+2b^2c^2 \] meaning
\[ \sqrt{\dfrac{2a^2b^2-a^4-b^4-c^4+2a^2c^2+2b^2c^2}{16}} \]
\[ =\sqrt{\dfrac{\left(a+b+c\right)\left(-a+b+c\right)\left(a-b+c\right)\left(a+b-c\right)}{16}} \]
\[ =\sqrt{\left(\frac{a+b+c}{2}\right)\left(\frac{-a+b+c}{2}\right)\left(\frac{a-b+c}{2}\right)\left(\frac{a+b-c}{2}\right)} \]
Now recall that $s=\frac{1}{2}(a+b+c)$. Thus
\[ \sqrt{\left(\frac{a+b+c}{2}\right)\left(\frac{-a+b+c}{2}\right)\left(\frac{a-b+c}{2}\right)\left(\frac{a+b-c}{2}\right)} \]
\[ =\sqrt{s(s-a)(s-b)(s-c)} \] therefore proving Heron's Formula.

\section*{1}
In a triangle with sides of length $a,b,$ and $c,\frac{1}{2}(a+b+c)$ is called the semi-perimeter, or half-perimeter.

\section*{3}
Law of cosines.

\section*{4}
Law of sines.

\section*{5}
Law of sines.

\section*{6}
Law of cosines.

\section*{7}
Law of sines.

\section*{10}


\end{document}
